\section{\method}
\label{sec:method}

Sections~\ref{sec:collapse}--\ref{sec:anchor} leave one objective standing. Assign each
scale $r\in\Rset$ its own readout temperature $\tau_r$, keep the zero-step member of the
family so that one target has full support, and add nothing else:
\begin{equation}
    \boxed{\;
    \mathcal{L}_{\method}
    \;=\;
    \frac{1}{|B|}\sum_{i\in B}\ \sum_{r\in\Rset}\ \omega_r\,
    \KL\!\left(p_{i,r}\ \middle\Vert\ q_i^{\tau_r}\right),
    \qquad
    \Rset=\{0\}\cup\mathcal{R},
    \;}
    \label{eq:objective}
\end{equation}
with $p_{i,r}$ from \eqref{eq:target}, $q_i^{\tau_r}$ from \eqref{eq:readout}, and
$\tau_r$ increasing over $r\in\mathcal{R}$ so that broader targets are read out more
smoothly. There is one term and no balancing weight: the quantity a pointwise anchor was
meant to supply is produced by the bank itself (Lemma~\ref{lem:gauge}), and the quantity a
separate relational term was meant to supply is the $r=0$ member of the family. Constants,
the choice of $\tau_0$, and the full procedure are in Appendix~\ref{app:impl}.

The scales share the underlying vector $v_i$ but not the distribution built from it. The
identity \eqref{eq:collapse} therefore no longer applies, and \eqref{eq:floor} becomes a
lower bound rather than an attained value --- so a run of \eqref{eq:objective} whose loss
falls below $\JS_\omega$ is direct evidence that the bank is doing something a single
readout cannot.

\paragraph{Shared columns.}
Anchors in a batch are scored against the union of all candidates in the batch,
$\mathcal{C}_B = \bigcup_{i\in B} C_i$, with $p_{i,r}(j) = 0$ for
$j \in \mathcal{C}_B\setminus C_i$ and $q_i^{\tau_r}$ masked at $j=i$. Its purpose is not
more negatives but \emph{shared} ones: Proposition~\ref{prop:levels}(b) removes the additive
freedom within an anchor, and shared columns extend the comparison across anchors, which is
the regime level-sensitive metrics measure. Because the candidate texts are encoded regardless,
this is free in forward cost; but it also multiplies the number of zero-target columns, so
it amplifies Proposition~\ref{prop:support} and is compatible with \eqref{eq:objective} only
because $r=0$ is present.

\paragraph{Sampling can reintroduce the collapse.}
Scoring an anchor against all $N-1$ columns is infeasible, so $C_i$ is drawn from some
$\mu_i$ and redrawn each epoch. The analysis requires one property of $\mu_i$.

\begin{condition}[Per-scale coverage]
\label{cond:cover}
$\Pr_{\mu_i}\!\left[\ \supp(p_{i,r}) \cap C_i \neq \emptyset\ \right] = 1$ for every
$r\in\Rset$.
\end{condition}

Condition~\ref{cond:cover} forbids the obvious sampler. Drawing the positive quota from the
mixture $\bar p_i = \sum_r \omega_r p_{i,r}$ is natural but wrong: its support is dominated
by the broadest walk, so the sharpest scale receives one or two scored candidates and its KL
term stops carrying ranking information --- reintroducing by sampling exactly the collapse
that Proposition~\ref{prop:collapse} identifies in the objective. A per-scale quota
satisfies the condition; Appendix~\ref{app:impl} gives the allocation we use.

\paragraph{The family is checkable before training.}
Each result of Section~\ref{sec:analysis} names a quantity that can be computed from the
teacher artifact alone, so a degenerate scale family is detectable before the first gradient
step:
\begin{equation}
    \underbrace{\KL\!\left(p_{i,r}\,\Vert\,\mathcal{U}_{\supp(p_{i,r})}\right)}_{\text{sharpness}},
    \qquad
    \underbrace{\KL\!\left(p_{i,r}\,\Vert\,p_{i,r'}\right)}_{\text{distinctness}},
    \qquad
    \underbrace{\JS_\omega\!\left(\{p_{i,r}\}\right)}_{\text{floor }\eqref{eq:floor}} .
    \label{eq:diag1}
\end{equation}
Sharpness near $0$ means the target is a binary neighbour label carrying no ranking signal;
distinctness near $0$ means $r$ and $r'$ are duplicates; a floor near $0$ certifies the
family is degenerate by Corollary~\ref{cor:floor}. These also select $\tau_g$ and bound a
useful $\mathcal{R}$ without a downstream sweep, and they matter during training as well:
by Corollary~\ref{cor:floor} the loss value alone separates neither a student that has
learned the neighbourhood ranking from one whose embeddings have collapsed towards uniform,
nor a converged student from an exhausted objective (Appendix~\ref{app:impl}).

\paragraph{Cost.}
Graph construction is a one-off $O(Nk)$ sparse operation over cached embeddings, and the
$r=0$ target is a lookup in the same cache. Training never touches the teacher: the memory
resident during optimisation is the student plus an $N\times d_T$ table of teacher vectors,
and the additional runtime cost over a mini-batch relational objective is $|\Rset|$ softmaxes
over $|\mathcal{C}_B|$ columns per anchor, negligible next to the encoder forward pass.
